\section{Formal details for realized processes and receiver presentations}
\label{app:process-presentation-details}

This appendix records formal properties used by the main-text typing of realized neural processes and receiver-level mechanism presentations. It also formalizes architecture-specific decoder classes as restrictions of the maximal receiver-level decoder functor: the main text's decoder hypothesis classes appear here as subfunctors. The main text keeps the constructions that drive the argument; categorical and almost-sure bookkeeping is collected here.

\subsection{State transport and the realized neural process}

Let $\SBor$ denote the category of standard Borel spaces and measurable deterministic maps. For a measurable map $f:X\to Y$, pushforward defines
\[
f_*:\Prob(X)\to\Prob(Y),
\qquad
(f_*\mu)(B)=\mu(f^{-1}(B)).
\]
Thus
\begin{equation}
\Prob:\SBor\longrightarrow\Set
\label{eq:app-probability-functor}
\end{equation}
is a covariant functor.

Fix represented cuts
\[
[L]=\{0<1<\cdots<L\}
\]
regarded as a category, and let
\[
X^\theta:[L]\to\SBor
\]
be a deterministic feed-forward neural process. Write
\[
F_{q\rightsquigarrow r}^\theta:=X^\theta(q\to r),
\qquad
A_q^\theta:=F_{0\rightsquigarrow q}^\theta.
\]
For a source state $\mu_0\in\Prob(X_0)$, set
\[
\mu_q:=(A_q^\theta)_*\mu_0.
\]
The category of elements $\int\Prob$ has objects $(X,\mu)$ and morphisms
\[
f:(X,\mu)\to(Y,\nu)
\]
with $f:X\to Y$ measurable and $f_*\mu=\nu$.

\begin{proposition}[The realized neural process]
\label{thm:app-state-decorated-process}
For every deterministic neural process $X^\theta$ and source state $\mu_0$, the pushed-forward states define a unique lift
\[
\widetilde X^{\theta,\mu_0}:[L]\longrightarrow\int\Prob
\]
whose projection to $\SBor$ is $X^\theta$ and whose value at cut $q$ is $(X_q,\mu_q)$.
\end{proposition}

\begin{proof}
For $q\le r$,
\[
(F_{q\rightsquigarrow r}^\theta)_*\mu_q
=(F_{q\rightsquigarrow r}^\theta)_*(A_q^\theta)_*\mu_0
=(A_r^\theta)_*\mu_0
=\mu_r.
\]
Hence every network arrow is a morphism in $\int\Prob$. The source state determines every downstream state through the unique arrow $0\to q$, so the lift is unique.
\end{proof}

\begin{proposition}[Equivalent states remain equivalent downstream]
\label{prop:app-measure-class-pushforward}
Let $\mu$ and $\nu$ be probability measures on $X$ with the same null measurable sets. For every measurable deterministic map $f:X\to Y$,
\[
f_*\mu\sim f_*\nu.
\]
Consequently, if two source states lie in the same measure class, their induced states lie in the same measure class at every represented cut of a deterministic neural process.
\end{proposition}

\begin{proof}
For every measurable $B\subseteq Y$,
\[
(f_*\mu)(B)=0
\Longleftrightarrow
\mu(f^{-1}(B))=0
\Longleftrightarrow
\nu(f^{-1}(B))=0
\Longleftrightarrow
(f_*\nu)(B)=0.
\]
Apply the result to each prefix map $A_q^\theta$.
\end{proof}

\subsection{Receiver refinement as a presentation category}

Fix a finite declared receiver family $J$ at one represented cut, with measurable maps
\[
Q_j:X\to Z_j.
\]
Write
\[
\Bool_J:=(2^J,\subseteq)
\]
for the Boolean receiver poset regarded as a category. For $K\subseteq J$, let
\[
Z_K:=\prod_{j\in K}Z_j,
\qquad
Q_K:=\langle Q_j\rangle_{j\in K},
\]
and for $K\subseteq K'$ let
\[
p_{K'K}:Z_{K'}\to Z_K
\]
be coordinate projection. Fix a standard Borel action space $\mathsf A$ and define
\[
\Dec_{\mathsf A}(K)
:=\{h:Z_K\to\mathsf A\mid h\text{ measurable}\}.
\]
Receiver refinement pulls a rule on $K$ back to one on $K'$,
\begin{equation}
\Dec_{\mathsf A}(K)\longrightarrow\Dec_{\mathsf A}(K'),
\qquad
h\longmapsto h\circ p_{K'K}.
\label{eq:app-decoder-refinement-map}
\end{equation}
These maps are functorial, so
\[
\Dec_{\mathsf A}:\Bool_J\to\Set
\]
is a functor. Its category of elements $\int\Dec_{\mathsf A}$ has objects $(K,h)$ with $K\subseteq J$ and $h:Z_K\to\mathsf A$.

A refinement arrow preserves the induced action pointwise. If $K\subseteq K'$ and $h'=h\circ p_{K'K}$, then
\begin{equation}
h'Q_{K'}=hQ_K.
\label{eq:app-refinement-preserves-action}
\end{equation}

\subsection{Architecture-specific decoder subfunctors}

The maximal decoder functor separates receiver access from any particular downstream implementation. Categorically, refinement-compatible decoder hypothesis classes are exactly subfunctors of the maximal decoder functor,
\[
\mathscr H\hookrightarrow\Dec_{\mathsf A}.
\]
Concretely, this consists of subsets
\[
\mathscr H_K\subseteq\Dec_{\mathsf A}(K)
\]
for each $K\subseteq J$ such that
\begin{equation}
h\in\mathscr H_K
\quad\Longrightarrow\quad
h\circ p_{K'K}\in\mathscr H_{K'}
\qquad(K\subseteq K').
\label{eq:app-architecture-refinement-closure}
\end{equation}
This is exactly the closure needed for passive receiver refinement to preserve the admissible downstream class.

For a probability state $\mu$ and measurable loss $\ell:Y\times\mathsf A\to[0,\infty]$, define
\[
\delta_{\mu,\mathscr H}(K)
:=
\inf_{h\in\mathscr H_K}
\mathbb E_\mu\!\left[\ell\bigl(\Phi(X),h(Q_K(X))\bigr)\right].
\]

\begin{proposition}[Refinement monotonicity for decoder subfunctors]
\label{prop:app-architecture-refinement-monotonicity}
If $\mathscr H\hookrightarrow\Dec_{\mathsf A}$ is a refinement-compatible decoder subfunctor and $K\subseteq K'$, then
\[
\delta_{\mu,\mathscr H}(K')\le\delta_{\mu,\mathscr H}(K).
\]
\end{proposition}

\begin{proof}
For $h\in\mathscr H_K$, closure \eqref{eq:app-architecture-refinement-closure} gives $h':=h\circ p_{K'K}\in\mathscr H_{K'}$. Equation~\eqref{eq:app-refinement-preserves-action} gives $h'Q_{K'}=hQ_K$ pointwise, so the two rules have the same risk. Taking the infimum over $h\in\mathscr H_K$ gives the claim.
\end{proof}

The maximal case is $\mathscr H=\Dec_{\mathsf A}$. Smaller subfunctors specialize the receiver-sufficiency problem to restricted downstream computation classes without changing the access maps themselves. A family that is not closed under \eqref{eq:app-architecture-refinement-closure} need not inherit the refinement monotonicity law.

In the finite exact setting, define
\[
\Rcal^{\mathscr H}_\Phi(S)
:=
\{K\subseteq J:\exists h\in\mathscr H_K\text{ with }\Phi|_S=hQ_K|_S\}.
\]
Then
\begin{equation}
\Rcal^{\mathscr H}_\Phi(S)\subseteq\Rcal_\Phi(S).
\label{eq:app-restricted-exact-subfamily}
\end{equation}
Hence failure of the maximal factorization condition is an access obstruction shared by every downstream architecture using only $Q_K(X)$, while membership in $\Rcal_\Phi(S)$ need not imply membership in $\Rcal^{\mathscr H}_\Phi(S)$ for a restricted architecture class.

\subsection{State-relative presentation equivalence}

For a probability state $\mu$ on $X$, define
\[
(K,h)\equiv_\mu(K',h')
\quad\Longleftrightarrow\quad
hQ_K=h'Q_{K'}
\qquad\mu\text{-almost surely}.
\]

\begin{proposition}[State-relative presentation equivalence]
\label{prop:app-state-relative-presentation-equivalence}
For every probability state $\mu$, the relation $\equiv_\mu$ is an equivalence relation on receiver-level presentations with common action space $\mathsf A$. Moreover:
\begin{enumerate}[label=(\roman*)]
\item every refinement arrow in $\int\Dec_{\mathsf A}$ relates equivalent presentations pointwise and therefore under every state;
\item if $\mu\sim\nu$ have the same null sets, then $\equiv_\mu$ and $\equiv_\nu$ coincide;
\item if $\bar\mu=\sum_{c=1}^m\lambda_c\mu_c$ with every $\lambda_c>0$, then
\[
(K,h)\equiv_{\bar\mu}(K',h')
\quad\Longleftrightarrow\quad
(K,h)\equiv_{\mu_c}(K',h')
\text{ for every }c.
\]
\end{enumerate}
\end{proposition}

\begin{proof}
Reflexivity, symmetry, and transitivity are inherited from almost-sure equality of measurable maps. Part (i) follows from \eqref{eq:app-refinement-preserves-action}. For part (ii), the disagreement set
\[
B:=\{x:h(Q_K(x))\neq h'(Q_{K'}(x))\}
\]
is measurable and is $\mu$-null if and only if it is $\nu$-null. For part (iii),
\[
\bar\mu(B)=\sum_{c=1}^m\lambda_c\mu_c(B).
\]
With all coefficients strictly positive, this vanishes exactly when every $\mu_c(B)$ vanishes.
\end{proof}

These properties justify the main text's hierarchy
\[
K\longrightarrow(K,h)\longrightarrow[hQ_K]_\mu
\]
and make architecture-specific specialization explicit without requiring the categorical details to interrupt the main argument.
